\chapter{Introduction}
\label{chap:introduction}

\section{Background}

Travel planning brings together information that is usually scattered across maps, business listings, booking sites, public transport pages, and local recommendations. A traveler may begin with a broad question such as where to eat near a landmark, then refine it with a location, a date, a preferred pace, or a limited amount of time. A useful assistant therefore needs to do more than produce fluent text. It must retrieve places that exist, preserve their geographic identity, make the limits of incomplete information visible, and present a schedule that can be understood and changed.

This problem is especially concrete in a destination catalogue assembled from community-maintained geographic data. OpenStreetMap is a collaborative source of geographic features, and its coverage and tagging vary by place and feature type \cite{haklay2008openstreetmap}. A record can have a name and coordinates while lacking opening hours, prices, or a street address. A system that fills those gaps with plausible-sounding guesses may appear helpful while making a traveler less informed. The design challenge is thus twofold: use structured records to ground place identity, and retain explicit uncertainty where records are incomplete.

Large language models (LLMs) have made it practical to interpret open-ended questions and produce natural-language responses. Transformer models provide a flexible foundation for this interaction \cite{vaswani2017attention,brown2020language}. However, a model's linguistic fluency does not guarantee that an attraction exists, that a price is current, or that a proposed trip can be completed on time. Retrieval-augmented generation (RAG) addresses part of this problem by conditioning generation on retrieved evidence \cite{lewis2020retrieval,gao2023retrieval}. It does not, by itself, make the data current or complete, prove that the retrieval is relevant, or validate a schedule.

The repository examined in this report, \textit{ChatbotTravel\_server} (also branded in the source documentation as \textit{VietTravel AI}), is a Python web application for travel search and itinerary management. It uses Flask and Flask-RESTX to expose HTTP endpoints, SQLAlchemy to model relational data, and PostgreSQL as its intended database. A local web client is served from the same Flask application. The data pipeline prepares tourism-related points of interest (POIs) from Vietnam OSM shapefiles and loads normalized records into a PostgreSQL staging table.

There are two related but distinct interaction paths in the codebase. The travel-chat path detects language, extracts intent and features, searches the staging data, and uses configured text-generation providers when available. The itinerary path accepts a structured planning request, retrieves candidate POIs, builds a draft with a deterministic scheduling routine by default, validates it, and returns it without persistence. A separate confirmation operation checks a bearer token and persists a confirmed schedule. Keeping these paths distinct is important: the source code does not implement one autonomous graph in which every chatbot request passes through the itinerary planner.

\section{Problem Statement}

The project addresses a practical engineering question:

\begin{quote}
How can a web-based travel assistant combine conversational access to destination data with a constrained itinerary workflow while making the source and certainty of place, schedule, and route information explicit?
\end{quote}

The question decomposes into several implementation concerns:
\begin{enumerate}[leftmargin=1.8em,itemsep=4pt]
    \item \textbf{Geographic grounding.} A generated answer should use POIs retrieved from the requested geographic scope and should retain the source identifier and coordinates needed to trace a recommendation back to the catalogue.
    \item \textbf{Incomplete attributes.} Community records do not consistently contain price, rating, contact information, or opening hours. Unknown values must remain unknown rather than being silently substituted with invented values.
    \item \textbf{Schedule validity.} A generated one-day schedule needs unique stops, valid coordinates, properly formatted time intervals, no overlaps, and a reasonable upper bound on the day span. Geographic distance and actual driving time are different quantities and should not be conflated.
    \item \textbf{Persistence and identity.} A plan assembled from OSM identifiers is not automatically a record in the application's curated \code{Attractions} table. The draft and confirmation stages must preserve this distinction and verify the authenticated user before writing.
    \item \textbf{Evaluability.} A useful capstone report should separate source-code behavior, dataset manifests, and test definitions from measured model quality or runtime performance. These evidence types answer different questions.
\end{enumerate}

\section{Objectives}

The engineering objectives reflected in the inspected source tree are:
\begin{enumerate}[leftmargin=1.8em,itemsep=2pt,parsep=0pt,topsep=3pt]
    \item Provide a REST API and browser interface for travel questions and itinerary planning.
    \item Prepare Vietnam OSM data as named POIs with standardized categories, coordinates, and source metadata.
    \item Search PostgreSQL by text or category, and by radius when coordinates are available.
    \item Build one-day itinerary drafts from retrieved candidates, using the deterministic planner by default.
    \item Validate each draft in application code and report errors separately from warnings.
    \item Let verified users confirm a draft and save it transactionally.
    \item State implementation limits and propose an evaluation plan without presenting unrun experiments as results.
\end{enumerate}

\section{Scope and Non-Goals}

The report covers the contents of the repository snapshot, including the Flask entry point, application factory, controllers, service modules, ORM models, SQL migrations, data-import scripts, OSM files, the pilot JSON fixture, and the current regression-test source. The report describes the intended local development setup from \code{docker-compose.dev.yml} and \code{requirements.txt}.

The work examined here does not establish that a particular PostgreSQL instance currently contains all records described in the data manifest. It does not provide an independently labeled relevance dataset, a human evaluation panel, a reproducible load-test log, or a model-quality benchmark. The source includes optional integration code for Google Gemini in the planner and a separate provider path in the chatbot; the itinerary controller explicitly calls the planner with external API execution disabled. Consequently, the default itinerary result is produced by a deterministic heuristic rather than by a live Gemini call. The codebase also contains optional ChromaDB-related code, but the itinerary orchestrator explicitly disables it when it searches PostgreSQL.

The project should therefore be interpreted as a working prototype with implementable interfaces and validation rules, not as evidence of production-scale coverage or measured recommendation quality. In particular, a static validator can reject malformed data but cannot guarantee a correct road route, current business hours, accessibility, price, safety, or user satisfaction.

\section{Method and Evidence Policy}

This report was prepared by inspecting the repository's source files and included documentation. Its implementation claims refer to named modules and executable paths. Its dataset counts refer to \url{data/osm_full_ingestion_summary.json} and are treated as manifest values rather than a live database count. Its testing discussion describes the cases in \url{tests/test_multi_agent_flow.py} and a local run on 21 October 2026 using the project virtual environment with Python 3.10.11. All nine cases passed; database and provider calls were mocked. The result applies to the local working tree at that time and is not a live integration, model-quality, or performance benchmark.

Evidence is classified in three ways throughout the report:
\begin{description}[leftmargin=2.1cm,style=standard,itemsep=1pt,parsep=0pt]
    \item[Recorded] A value appears in a committed data file, configuration file, migration, or test definition.
    \item[Implemented] A behavior is traceable to a function or controller path in the source code.
    \item[Unmeasured] A performance, quality, or operational outcome for which the snapshot provides no reproducible run output.
\end{description}
This policy is particularly relevant to Chapter~\ref{chap:results_discussion}: the manifest provides descriptive counts, and test source provides test intentions, but neither is equivalent to an executed end-to-end benchmark.

\section{Contributions and Report Organization}

The repository's principal engineering contribution is an end-to-end structure connecting geographic data preparation to a web-facing travel application. The data loader filters and normalizes selected OSM layers; a relational staging model preserves the source identifier and source attributes; search services provide database-backed candidates; the planner produces a time-ordered draft; and a validator checks a defined set of structural and scheduling constraints. Authentication and itinerary persistence are implemented as complementary application features.

Chapter~\ref{chap:related_work} reviews recommendation, grounding, agent orchestration, and geographic planning concepts relevant to the implemented design. Chapter~\ref{chap:dataset} describes the available OSM evidence, import rules, schema, and data-quality caveats. Chapter~\ref{chap:model_configuration} explains the running architecture and its model/provider configuration. Chapter~\ref{chap:results_discussion} reports only reproducible evidence found in the repository, discusses limitations, and proposes a future evaluation protocol. Appendices provide a compact data dictionary, endpoint inventory, and the test and evaluation checklist.
